% Generated draft: requires mathematical and source review.
\documentclass[11pt]{article}
\usepackage[margin=25mm]{geometry}
\usepackage{fontspec}
\setmainfont{Latin Modern Roman}
\usepackage{amsmath,amssymb,mathtools,unicode-math}
\setmathfont{Latin Modern Math}
\usepackage{enumitem,needspace,xurl,hyperref}
\setlength{\parindent}{0pt}
\setlength{\parskip}{5pt}
\setlength{\emergencystretch}{3em}
\newcommand{\N}{\mathbb{N}}
\newcommand{\Z}{\mathbb{Z}}
\newcommand{\Q}{\mathbb{Q}}
\newcommand{\R}{\mathbb{R}}
\newcommand{\C}{\mathbb{C}}
\newcommand{\F}{\mathbb{F}}
\newcommand{\A}{\mathbb{A}}
\newcommand{\PP}{\mathbb{P}}
\newcommand{\E}{\mathbb{E}}
\newcommand{\eps}{\varepsilon}
\newcommand{\dd}{\,\mathrm{d}}
\begin{document}
\section*{UnsolvedMath ID: 30002791}
\textbf{Source problem code:} catalogue entry
\subsection{Definitions and mathematical statement}
The notation \(H(C)\) denotes the Harbourne constant referred to in the supplied record. The record does not define this invariant, its domain, or its conventions, so no formula for \(H(C)\) is imposed here.

The recorded existence question is
\[
\text{Does there exist an irreducible curve }C\text{ with }H(C)\leq -2?
\]
The accompanying background discusses plane curves and distinguishes positive characteristic from the complex numbers. It therefore suggests the intended unresolved specialization
\[
\text{Does there exist an irreducible plane curve }C\text{ over }\C
\text{ with }H(C)\leq -2?
\]
This specialization is inferred from the supplied background; the recorded question itself does not specify the ground field or ambient surface. A complete mathematical formulation requires the intended definition of \(H(C)\), the admissible class of curves, and any hypotheses on singularities or other boundary cases.

The record also asks for analogous questions or constructions on other surfaces. It does not specify a class of surfaces, an invariant on those surfaces, or a quantified existence or boundedness assertion. That part is retained as an unresolved research direction rather than a definite mathematical proposition.
\subsection{Short English statement}
The record asks whether an irreducible curve can have Harbourne constant at most \(-2\), with the background pointing to complex plane curves as the unresolved case, and asks for analogues on other surfaces. The invariant and the precise hypotheses are missing, so this is a qualified statement-only draft.
\subsection{Sources}
\begin{itemize}
\item Supplied local UnsolvedMath registry record, permanent ID \(30002791\), source problem code \texttt{catalogue entry}, entitled ``Irreducible Curves with Strongly Negative Harbourne Constants.'' The proposed year is 2015; no individual proposer is identified.
\item Supplied workshop attribution: \emph{Mini-Workshop: Ideals of Linear Subspaces, Their Symbolic Powers and Waring Problems} (2015), \emph{Oberwolfach Reports}, pp.~489--532, DOI 10.4171/owr/2015/9. Original report: \url{https://doi.org/10.4171/owr/2015/9}. The supplied background describes these reports as documenting mathematical workshops and their open-problem discussions.
\item Supplied literature reference, labeled ``Dimca, Harbourne, and Sticlaru, bounded negativity and singular curves'': \url{https://arxiv.org/abs/2101.07187}.
\item Literature assessment dated exactly \texttt{2026-08-21}: the supplied status is \texttt{partially\_solved}. The supplied assessment says that irreducible plane curves with \(H(C)\leq -2\) are known in positive characteristic, that strong lower bounds hold in several restricted cases, and that existence over the complex numbers and the analogous surface-wide questions remain unresolved. This is supplied dated curation derived from the OWR record, not an independently verified literature assessment; its status and cited claims require independent verification.
\item The registry labels the statement \texttt{corrected\_verified} and says that mathematical content was preserved while notation, typography, or omitted definitions were repaired for publication. This supplied editorial label does not resolve the missing definition or establish mathematical correctness here.
\end{itemize}
\end{document}
